\section{Results and Synthesis}

This section consolidates the findings of the 37 included studies. The synthesis follows a hybrid thematic approach: six themes were derived inductively from the structured fields of the extraction records (artifacts, roles, execution, validation, portability, risks, nature, and evidence) and then confronted with the six-dimension taxonomy proposed in the previous section, so as to test it against the evidence rather than assume it. The themes answer RQ1 through RQ4; RQ5, on gaps, is addressed in the Research Agenda. The strength of each study's evidence is weighted by a source-category quality assessment, described in the Method and summarized below (Table~\ref{tab:qualidade}).

\noindent\textbf{Quantitative overview.} Table~\ref{tab:quant} summarizes the composition of the corpus. The framework category dominates almost entirely (34 of the 37 studies), which empirically supports the thesis that the unit of work has shifted from the isolated prompt to the structured process. Empirical evaluation is claimed by 26 studies and 22 release a public repository, but ten are arXiv preprints not yet peer-reviewed, which calls for a qualified reading of the strength of the evidence (Figure~\ref{fig:evidence}). The temporal distribution shows steady acceleration of the field, from a single study in 2022 to eleven in 2026, with most of the corpus concentrated from 2024 onward (Figure~\ref{fig:year}). Venues are dispersed, with the ACM Transactions on Software Engineering and Methodology as a recurring outlet (five studies) followed by a long tail of conferences and journals with one or two studies each, consistent with an emerging field that still lacks a consolidated forum.

\begin{table}[h!]
\centering
\caption{Composition of the corpus of 37 included studies. Counts extracted from the extraction matrix.}
\label{tab:quant}
\small
\begin{tabularx}{\textwidth}{@{}p{0.50\textwidth}r Y@{}}
\toprule
\textbf{Aspect} & \textbf{N} & \textbf{Note} \\
\midrule
Total included studies & 37 & After two screening passes and three snowballing rounds \\
Nature: framework & 34 & Near-total predominance \\
Nature: architecture / platform / method & 1 / 1 / 1 & JulesWhite2023, YuCheng2023, JSauvola2024 \\
With empirical evaluation & 26 & Some degree of validation claimed \\
With public repository & 22 & Code traceability in more than half \\
arXiv preprint & 10 & Literature not yet peer-reviewed \\
Survey or review & 2 & AhmedEHassan2024, JundaHe2025 \\
Case study & 5 & Ulfsnes2024, ZeeshanRasheed2024, Baron2025, Khan2026, Macedo2026 \\
Distribution by year & \multicolumn{2}{l}{2022:1, 2023:4, 2024:10, 2025:11, 2026:11} \\
Academic or technical / industrial origin & 33 / 4 & Four studies from industry or professional practice \\
High extraction confidence & 35 & Medium in Dam2025 and JundaHe2025 \\
\bottomrule
\end{tabularx}
\end{table}

\begin{figure}[h!]
\centering
\includegraphics[width=0.70\textwidth]{studies_by_year_en.pdf}
\caption{Studies by publication year. The field accelerates from a single study in 2022 to a plateau of ten to eleven per year from 2024 onward.}
\label{fig:year}
\end{figure}

\begin{figure}[h!]
\centering
\includegraphics[width=0.80\textwidth]{evidence_profile_en.pdf}
\caption{Evidence profile of the corpus. Most studies claim empirical evaluation and release a public repository, but ten are arXiv preprints not yet peer-reviewed.}
\label{fig:evidence}
\end{figure}

\noindent\textbf{Quality assessment by source category.} Table~\ref{tab:qualidade} summarizes the quality assessment described in the Method, with the four-item checklist applied to the 37 studies. Eleven studies (30\%) fall in the high-quality band, seventeen (46\%) in the medium band, and nine (24\%) in the low band. Most of the corpus comes from peer-reviewed venues (25 studies, 68\%), against ten arXiv preprints (27\%) and two non-indexed technical records (5\%). At the item level, 26 studies (70\%) report empirical evaluation and 22 (59\%) release a public repository, but only five (14\%) bring applied validation in a case study or real project, which bounds how far the current evidence supports production use. The low band concentrates preprints and technical records without recorded empirical evaluation; their findings are read cautiously in the synthesis and no theme rests exclusively on them. The assessment is conservative by construction, since it derives from the signals recorded in the extraction rather than from a full-text reappraisal of each study.

\begin{table}[h!]
\centering
\caption{Source-category quality assessment of the 37 studies. Items derived from observable signals in the extraction matrix; QA is the sum of QA1 to QA4 (0 to 4).}
\label{tab:qualidade}
\small
\begin{tabularx}{\textwidth}{@{}p{0.46\textwidth}r Y@{}}
\toprule
\textbf{Category} & \textbf{N} & \textbf{\%} \\
\midrule
High quality (QA $\geq$ 3) & 11 & 30\% \\
Medium quality (QA = 2) & 17 & 46\% \\
Low quality (QA $\leq$ 1) & 9 & 24\% \\
\midrule
Peer-reviewed venue (QA1) & 25 & 68\% \\
arXiv preprint & 10 & 27\% \\
Non-indexed technical record & 2 & 5\% \\
\midrule
Reported empirical evaluation (QA2) & 26 & 70\% \\
Applied validation in a case study (QA3) & 5 & 14\% \\
Public artifact for reproducibility (QA4) & 22 & 59\% \\
\bottomrule
\end{tabularx}
\end{table}

\noindent\textbf{Six inductive themes.} The analysis surfaced six themes that, together, cover the 37 studies (100\% coverage, above the 80\% target). Each study may belong to more than one theme; the order reflects decreasing frequency.

\noindent\textbf{T1: multi-agent orchestration and role division.} This is the most populous theme and the dominant architectural axis of the corpus, answering RQ1 directly. The central convergence is that structuring the work by distinct roles, not the choice of language model, is what differentiates the frameworks. ChatDev encodes a chat chain with design, coding, testing, and review phases, plus a communicative dehallucination mechanism \cite{ChenQian2023}. MetaGPT internalizes standard operating procedures from software engineering, assigning specific artifacts to each role: the product manager generates requirements, the architect generates the design and interfaces, the developer implements \cite{SiruiHong2023}. CodePori scales that arrangement to agents for management, architecture, structure, development, verification, and finalization \cite{ZeeshanRasheed2024}, and AutoDev and AutoCodeRover bring agent coordination to real program-improvement tasks \cite{MicheleTufano2024,YuntongZhang2024}. He et al.'s reviews and Hassan et al.'s Software Engineering 3.0 vision confirm that the literature already organizes these systems by lifecycle stage \cite{JundaHe2025,AhmedEHassan2024}. Consensus, however, is partial: there is clear tension over the degree of autonomy. Sapkota et al. explicitly contrast vibe coding and agentic coding \cite{RanjanSapkota2025}, and while CodePori and MetaGPT lean toward end-to-end autonomous development, other designs preserve human decision points. Reliability and coordination cost appear as recurring risks in this theme \cite{Dam2025,Alsegier2026,Merchant2025,Paduraru2026,SKaruppuchamy2026,Paiva2026,Akbar2025,LiyiCai2025,Anon2026,BESLEAGA2026}.

\noindent\textbf{T2: specification and contract as the central artifact.} This theme answers RQ1 and RQ2 and gives substance to the thesis that the specification becomes the unit of coordination. Spec-driven development formalizes the move from code to contract, treating the specification as a versionable artifact \cite{DeepakBabuPiskala2026}. AgileGen uses acceptance criteria as an executable contract and keeps memory and iteration until validation \cite{SaiZhang2025}, and MetaGPT already made structured requirements the source of truth across roles \cite{SiruiHong2023}. There is, however, a productive contradiction in the direction of the flow: most assume spec-first, where the specification precedes and generates the code, whereas Reversa inverts the flow, recovering operational specifications from legacy code to feed the agents \cite{Macedo2026}. This opposition between specifying first and recovering the specification from legacy is a relevant axis of differentiation. Risk-aware requirements engineering ties the specification to governance \cite{Watfa2026}, and the misalignment between specification and code is discussed as a risk across several studies \cite{DeepakBabuPiskala2026,Paduraru2026,Watfa2026,YuCheng2023,Gupta2022,MitulModi2024,Khan2026,Gadamsetty2025}.

\noindent\textbf{T3: context engineering and repository knowledge retrieval.} This theme answers RQ2 with a focus on the context dimension. Context engineering is treated as a discipline of its own for agents in open-source software \cite{SeyedmoeinMohsenimofidi2025}. AutoCodeRover shows that structural retrieval of context from the code, and not just textual search by embeddings, improves patch generation \cite{YuntongZhang2024}, while the synchronization of multimodal artifacts seeks to keep coherence across heterogeneous artifacts over the cycle \cite{Paiva2026}. Self-evolving software pushes the theme to its extreme \cite{LiyiCai2025}, but it sits in tension with the emphasis on traceability and contract in T2 and T4: the more the software modifies itself autonomously, the harder it is to anchor each change in an auditable specification \cite{Macedo2026,AhmedEHassan2024,Chechik2026}. The risk tied to context is the most frequent in the whole corpus, which suggests that the context window and context fidelity are the transversal practical bottleneck of the field.

\noindent\textbf{T4: validation, traceability, assurance, and governance.} This theme answers RQ2, in the validation dimension, and RQ4, in the risks of governance and traceability. Almost the entire corpus declares some validation mechanism, but here we group the studies for which validation is the central point rather than an appendix. Trace-based assurance articulates contracts, tests, and governance \cite{Paduraru2026}; verifiable collaboration among agentic assistants records the acts of the flow in an auditable way \cite{Merchant2025}; and product-line engineering is imported to govern the variability of agency \cite{Alsegier2026}. Traceability and trace auditing appear as a recurring requirement for trusting agents on real projects \cite{Paduraru2026,Merchant2025,Macedo2026,SKaruppuchamy2026,Watfa2026,Chechik2026,Anon2026}. There is, however, a cost contradiction: on one side, heavy governance ceremony \cite{Merchant2025,Paduraru2026}; on the other, lightweight agility aimed at startups and lean flows \cite{Khan2026,RanjanSapkota2025,JulesWhite2023}. Coordination cost is named as a risk in much of this group \cite{MitulModi2024,Ulfsnes2024,Akbar2025,Baron2025,KRRaghi2024,Gupta2022,KumarAle2024}, which materializes the tension between governance and productivity.

\noindent\textbf{T5: human-AI collaboration, adoption, and trust (emergent).} These studies focus not on the technical architecture but on integrating AI into human work: collaboration flows, organizational adoption factors, and trust building. A human-AI collaboration and adaptation framework models workflow factors and organizational strategies \cite{DanielRusso2024}; field evidence describes how real teams reorganize collaboration and flow \cite{Ulfsnes2024,Akbar2025}; and trust is treated as a precondition for adoption \cite{Baron2025}, with future scenarios of AI-assisted operation also documented \cite{JSauvola2024}. The notable finding is that this theme does not fit cleanly into any of the six taxonomy dimensions, all of which are technical and architectural: the human, organizational, and trust dimension cuts across the corpus but has no home in the current taxonomy. Six studies (16\% of the corpus) have this as their primary focus, enough to recommend extending the taxonomy.

\noindent\textbf{T6: security, robustness, and red-teaming (emergent).} These studies focus on the security of the code and of agentic behavior. ASTRA proposes autonomous spatial-temporal red-teaming for AI software assistants, treating security as an active activity rather than an assumed property \cite{XiangzheXu2025}; risk-aware requirements engineering and the discussion of the security implications of agentic coding reinforce the theme \cite{Watfa2026,RanjanSapkota2025,Merchant2025,BESLEAGA2026}, and ChatDev's communicative dehallucination touches the point from the reliability side \cite{ChenQian2023}. Security is among the most frequently cited risks, named in 24 of the 37 studies, behind context, reliability, and coordination cost (Figure~\ref{fig:risk}). Like T5, this theme connects only partially to the six dimensions: security appears almost always as a risk to mitigate, rarely as a positive design dimension embedded in the architecture. ASTRA is the exception that makes security the central object, which suggests treating it as a transversal evaluation criterion rather than as one of the dimensions that classify frameworks.

\begin{figure}[h!]
\centering
\includegraphics[width=0.80\textwidth]{risk_frequency_en.pdf}
\caption{Frequency of reported risks across the corpus. Context, reliability, and coordination cost lead; security, though prominent, ranks below them, and explicit lock-in tagging is rare.}
\label{fig:risk}
\end{figure}

\noindent\textbf{Confronting the taxonomy with the evidence.} Table~\ref{tab:tema-dim} crosses each inductive theme with the six taxonomy dimensions, indicating the strength of the relation, which is coded by the author from the extraction records. The reading is direct. Five of the six dimensions receive strong support from at least one theme: roles and execution anchor in T1; specification and validation, in T2 and T4; context, in T3. In that part, the taxonomy holds up against the evidence. Figure~\ref{fig:heatmap} renders the same mapping as a heatmap. The exception is portability, which receives no strong relation from any theme and is absent in two. In the data, portability appears almost always on the negative side, as a risk of lock-in and of dependence on an IDE, a commercial model, or a platform. The portability field of the records notes dependence on an IDE or environment in nearly all 37 studies, and lock-in is named explicitly as a risk in only one of them \cite{SeyedmoeinMohsenimofidi2025}. Portability is therefore the weakest dimension of the taxonomy when confronted with the evidence, and should be read as a predominantly diagnostic dimension, not as a positive differentiator between frameworks. Themes T5 and T6, in turn, have no clear home in the six dimensions: the human-organizational axis and the security axis emerge from the data as concrete recommendations for extension, and the rule of the hybrid mode was respected by not forcing them into the existing dimensions.

\begin{table}[h!]
\centering
\caption{Hybrid mapping: inductive themes against the six taxonomy dimensions. Legend: strong (the dimension is the theme's axis), medium (relevant presence), weak (marginal touch), absent (no significant relation).}
\label{tab:tema-dim}
\small
\begin{tabularx}{\textwidth}{@{}p{3.5cm}YYYYYY@{}}
\toprule
\textbf{Theme} & \textbf{Spec.} & \textbf{Context} & \textbf{Roles} & \textbf{Exec.} & \textbf{Valid.} & \textbf{Portab.} \\
\midrule
T1 Multi-agent orchestration & medium & medium & strong & strong & medium & weak \\
T2 Specification and contract & strong & medium & weak & medium & strong & weak \\
T3 Context engineering & weak & strong & weak & medium & weak & weak \\
T4 Validation and governance & medium & weak & weak & medium & strong & weak \\
T5 Human-AI collaboration & weak & weak & weak & weak & weak & absent \\
T6 Security and red-teaming & weak & weak & weak & medium & medium & absent \\
\bottomrule
\end{tabularx}
\end{table}

\begin{figure}[h!]
\centering
\includegraphics[width=0.92\textwidth]{theme_dimension_en.pdf}
\caption{Hybrid mapping of the six inductive themes against the six taxonomy dimensions. Five dimensions are anchored by at least one theme; portability stays weak throughout, and themes T5 and T6 fall partly outside the taxonomy.}
\label{fig:heatmap}
\end{figure}

\noindent\textbf{Coverage of the research questions.} Table~\ref{tab:tema-rq} crosses the themes with the research questions. RQ1 and RQ2 are answered by most themes; RQ3 and RQ4 cut across the whole corpus, and RQ5 is a meta-question answered transversally in the Research Agenda rather than by study-by-study tagging, which is why the individual records mark only RQ1 through RQ4.

\begin{table}[h!]
\centering
\caption{Inductive themes against the research questions.}
\label{tab:tema-rq}
\small
\begin{tabularx}{\textwidth}{@{}p{3.5cm}YYYYY@{}}
\toprule
\textbf{Theme} & \textbf{RQ1} & \textbf{RQ2} & \textbf{RQ3} & \textbf{RQ4} & \textbf{RQ5} \\
\midrule
T1 Multi-agent orchestration & yes & yes & yes & yes & via synthesis \\
T2 Specification and contract & yes & yes & partial & yes & via synthesis \\
T3 Context engineering & yes & yes & partial & yes & via synthesis \\
T4 Validation and governance & partial & yes & yes & yes & via synthesis \\
T5 Human-AI collaboration & partial & yes & yes & yes & via synthesis \\
T6 Security and red-teaming & no & partial & yes & yes & via synthesis \\
\bottomrule
\end{tabularx}
\end{table}
